% -*- TeX:RNW:UK -*-
\documentclass[10pt]{beamer}
\usetheme{metropolis}
%\useinnertheme{rectangles}
\setbeamercovered{%
still covered={\opaqueness<1->{15}},
again covered={\opaqueness<1->{40}}}

\hypersetup{colorlinks,linkcolor=black,urlcolor=brown,citecolor=brown}

\usepackage{amsmath,amssymb,amsthm}
\usepackage{unicode-math}

% We set the Lucida OTF fonts as default
\usepackage{fontspec}
\setmainfont{Lucida Bright OT}
\setsansfont{Lucida Sans OT}
\setmonofont{Lucida Console DK}[Scale=MatchLowercase]

\newfontfamily\typiconsfont{Typicons}[Scale=0.9]
\newcommand\typicons[1]{{\typiconsfont\symbol{#1}}}
\newcommand\Attention{\colorbox{yellow}{\textcolor{black}{\typicons{"E136}}}\xspace}
\newcommand\Advanced{\colorbox{black}{\textcolor{white}{\typicons{"E04E}}}\xspace}

\newfontfamily\webglyphsfont{WebHostingHub-Glyphs}[Scale=0.7]
\newcommand\webglyphs[1]{{\webglyphsfont\symbol{#1}}}
\newcommand\Discussion{\colorbox{white}{\textcolor{black}{\webglyphs{"F134}}}\xspace}
\newcommand\DExamples{\colorbox{black}{\textcolor{white}{\webglyphs{"F134} examples?}}}
\newcommand\Reading{\colorbox{black}{\textcolor{white}{\webglyphs{"F0C1}}}\xspace}

\newfontfamily\meteoconsfont{Meteocons}
\newcommand\meteocons[1]{{\meteoconsfont\symbol{#1}}}
\newcommand\meteosun{\meteocons{"0042}}
\newcommand\meteosuncloud{\meteocons{"0048}}
\newcommand\meteorain{\meteocons{"0052}}
\newcommand\meteosolidsun{\meteocons{"0031}}

\newfontfamily\KRfarmfont{KR Back On The Farm}
\newcommand\KRfarm[1]{{\KRfarmfont\symbol{#1}}}
\newcommand\farmplant{\KRfarm{"0049}}

\usepackage{polyglossia}
\setdefaultlanguage[variant = british, ordinalmonthday = false]{english}

\usepackage[style=authoryear-comp,firstinits,sortcites,maxcitenames=2,%
    mincitenames=1,maxbibnames=10,minbibnames=10,uniquename=mininit,%
    uniquelist=minyear,sortfirstinits=true]{biblatex}
\addbibresource{../references/ecophys.bib}
%\renewcommand{\bibfont}{\footnotesize}

\usepackage{abbrev}
\DeclareGraphicsExtensions{.png,.jpg,.jpeg,.pdf}

\begin{document}

\title{PBIO-141\\Sensory and Physiological Ecology of  Plants}
\subtitle{12: Carbon Assimilation (C)}
\author{Pedro J. Aphalo}
\date{January-February 2018}
\institute[Univ.\ of Helsinki]{M.Sc.\ in Plant Biology, University of Helsinki\\[2ex] \url{http://blogs.helsinki.fi/aphalo/}}


  \begin{frame}
    \maketitle
  \end{frame}

  \begin{frame}[c]
    \begin{center}
      \begin{small}
        \copyright 2006-2018 by Pedro J. Aphalo\\
       University of Helsinki, Finland.\\
        \textcolor{blue}{\url{http://blogs.helsinki.fi/aphalo/}}\\[2ex]
      \end{small}

      \begin{footnotesize}
        Sensory and Physiological Ecology of  Plants: `11: Carbon Assimilation (C)' by Pedro J. Aphalo is licensed under a Creative Commons Attribution-ShareAlike 4.0 International License.

        Illustrations and text quoted from copyrighted sources is excluded from this license and their use should respect the original licenses.\\[2ex]
      \end{footnotesize}

      \includegraphics[width=6em]{../figures/copyright/by-sa}
    \end{center}
  \end{frame}

  \begin{frame}
    \frametitle{Outline}
    \tableofcontents
  \end{frame}


\section{Response curves}

\begin{frame}{Light response curve}
    \begin{itemize}
        \item A light response curve is a plot of $A$
        (net photosynthetic rate = \emph{a}ssimilation rate) against
        $I$ (irradiance).%
        \item $A$ is usually expressed in moles of C per unit leaf
        area and unit time (\umol, for example).%
        \item $I$ is usually incident PAR photon irradiance (\umol) also called PPFD.
        \item less frequently absorbed light (\umol) is used instead of irradiance.
    \end{itemize}
\end{frame}

\begin{frame}{Photosynthesis limitation}
    \centering
    \includegraphics[height=0.8\textheight]{figures/Taiz9.7LightCO2LimitationRegions.eps}\\
    {\footnotesize Limitation of carbon assimilation by \COtwo or light
    supply
    \autocite[from][]{TaiZei2006}.}
\end{frame}

\begin{frame}{\COtwo response curve}
    \begin{itemize}
        \item A \COtwo response curve is a plot of $A$
        (photosynthetic rate = \emph{a}ssimilation rate) against
        $C_a$ (ambient \COtwo concentration) or $C_i$ (intercellular
        \COtwo concentration, the concentration in the air space of the leaf).
        \item $A$ is usually expressed in moles of C per unit leaf
        area and unit time (\umol, for example).%
        \item $C_a$ or $C_i$ is usually expressed as a molar fraction (\ppm), or less
        frequently as volume fraction (ppm) or partial pressure (Pa).
    \end{itemize}
\end{frame}

\begin{frame}{\COtwo response curve}
    \centering
    \includegraphics[height=0.75\textheight]{figures/Taiz9.22CO2ResponseC3C4.eps}\\
    {\footnotesize \COtwo response of $A$ in C3 and C4 plants
    \autocite[from][]{TaiZei2006}.}
\end{frame}

\section{Response to temperature}

\begin{frame}{Adaptation and acclimation}
    \begin{itemize}
        \item Under equal growing conditions, different species have
        different optimal temperatures for photosynthesis
        (adaptation).
        \item The optimal temperature for photosynthesis on the same
        plant can shift with seasons or growing conditions ---i.e. different ambient
        temperatures (acclimation).
        \item The capacity to acclimate may be different in
        different species, or in different populations of the same
        species.
    \end{itemize}
\end{frame}

\begin{frame}{Temperature optima}
    \centering
    \includegraphics[width=\textwidth]{figures/Crawley1.9PhotosynTemperatureSpecies.eps}\\
    {\footnotesize Adaptation of temperature optima \autocite[from][]{Crawley1997}.}
\end{frame}

\begin{frame}{Temperature optima}
    \centering
   \includegraphics[width=\textwidth]{figures/Crawley1.10PhotosynTemperatureSeasons.eps}\\
    {\footnotesize Acclimation of temperature optima \autocite[from][]{Crawley1997}.}
\end{frame}

\begin{frame}{A vs.\ temperature}
    \centering
    \includegraphics[width=\textwidth]{figures/Lambers36PhotosynTemperature.eps}\\
    {\footnotesize Acclimation and adaptation of photosynthesis optimal
    temperature to growing temperature
    \autocite[from][]{LambersEtAl1998}.}
\end{frame}

\section{Water stress}

\begin{frame}{Water stress}
    \begin{itemize}
        \item When water is in short supply in the soil, or when there is low
        humidity in the air, stomata close, $C_i$ decreases, and $A$ decreases.
        \item If water stress persists for several days the
        ``photosynthetic machinery'' gets affected and even if we
        artificially increase $C_i$ by increasing $C_a$, $A$ does
        not return to the rate it had before the stress.
        \item If we water the plant, there may be an after effect
        lasting up to a few days. This effect is seen as partially
        closed stomata and lower $A$ than before the stress started.
        \item[-] (We will return to the effects of water stress when we
        study water relations.)
    \end{itemize}
\end{frame}

\begin{frame}{A vs.\ C$_i$}
    \centering
    \includegraphics[height=0.8\textheight]{figures/Lambers29PhotosynCi.eps}\\
    {\footnotesize \COtwo-response curves of photosynthesis for plants under drought
    \autocite[from][]{LambersEtAl1998}.}
\end{frame}

\section{Acclimation and adaptation to shade}

\begin{frame}{Sun and shade}
    \begin{itemize}
        \item There are shade and sun adapted species.
        \item In many species plants acclimate to the growing light
        environment.
        \item Some tree species produce very distinct sun and shade
        leaves on different part of the crown.
        \item When grown in controlled environments we call them
        high light (HL) and low light (LL) plants. (same light quality.)
        \item Natural shadelight differs in quality (colour) from
        sunlight, and also in irradiance (quantity).
    \end{itemize}
\end{frame}

\begin{frame}{Adaptation vs. acclimation to shade \Discussion 5 + 2 min}
    \centering
    \includegraphics[height=0.75\textheight]{figures/Taiz9.8LightResponseSunShadeSpecies.eps}
    \includegraphics[height=0.75\textheight]{figures/Taiz9.10LightResponseSunShadeGrowth.eps}\\
    {\footnotesize Adaptation and acclimation to shade in different species (look also at the start of the curves
    at low irradiances).
    \autocite[from][]{TaiZei2006}.}
\end{frame}

%\begin{frame}{Acclimation to low light}
%    \centering
%    \includegraphics[height=0.75\textheight]{figures/Taiz9.10LightResponseSunShadeGrowth.eps}\\
%    {\footnotesize Acclimation happens in response to growth conditions
%    \autocite[from][]{TaiZei2006}.}
%\end{frame}
%
%\begin{frame}{Growth irradiance}
%    \centering
%    \includegraphics[width=\textwidth]{figures/Crawley1.5HLLLAtriplexPhotosyn.eps}\\
%    {\footnotesize Acclimation to growth irradiance \autocite[from][]{Crawley1997}.}
%\end{frame}
%
\begin{frame}{Acclimation to light \Discussion 5 + 5 min}
    What understanding do we get from these three plots?\\
    \centering
    \includegraphics[width=0.9\textwidth]{figures/Givnish1988-2.eps}\\
    {\footnotesize Photosynthesis expressed on different bases: area, mass, or protein \autocite[from][]{Givnish1988}.}
\end{frame}

\begin{frame}{Growth irradiance: morphology and anatomy}
    \centering
    \includegraphics[height=0.6\textheight]{figures/HLLLplants-1.eps}\ 
    \includegraphics[height=0.4\textheight]{figures/HLLLleaves-1.eps}\\
    {\footnotesize Acclimation to growth irradiance in radish
    (\species{Raphanus sativus})
    \autocite[from][]{Lichtenthaler1985}.}
\end{frame}

\begin{frame}{Growth irradiance}
    \centering
    \includegraphics[height=0.85\textheight]{figures/HLLLphotosynFuchsia-1.eps}\\
    {\footnotesize Acclimation to growth irradiance in \species{Fuchsia
    magellanica} (one clone)
    \autocite[from][]{AphaloEtAl1991}.}
\end{frame}

\begin{frame}{}
    \centering
    \includegraphics[width=0.75\textwidth]{figures/Fuchsia_HL_photo.eps}\\
    \includegraphics[width=0.75\textwidth]{figures/Fuchsia_LL_photo.eps}
\end{frame}

\begin{frame}{Sun and shade leaves}
    \includegraphics[width=0.55\textwidth]{figures/SunShadeLeaves-1.eps}
    \includegraphics[height=0.4\textheight]{figures/SunShadeLeaves-2.eps}\\
    {\footnotesize Sun and shade leaves from different parts
    of the crown of the same beech (\species{Fagus sylvatica}) tree
    \autocite[from][]{Lichtenthaler1985}.}
\end{frame}

\begin{frame}{Sun and shade leaves}
    \centering
   \includegraphics[width=0.9\textwidth]{figures/SunShadeLeaves-3.eps}\\
   {\footnotesize Shape can also be different in sun and shade leaves, e.g. in oak \autocite[from][]{Purves1995}.}
\end{frame}


\begin{frame}{\species{Thermopsis}}
    \centering
    \includegraphics[height=0.70\textheight]{figures/ThermopsisPhoto-1.eps}\\
    {\footnotesize False lupin plants (\species{Thermopsis})}
    {\tiny (from \url{http://linnaeus.nrm.se/flora/di/faba/therm/thermon.html}).}
\end{frame}

\begin{frame}{Sun and shade plants}
    \centering
    \includegraphics[width=\textwidth]{figures/SunShadeLeafThermopsis-1.eps}\\
    {\footnotesize Acclimation of high and low light in \species{Thermopsis} \autocite[from][]{TaiZei2002}.}
\end{frame}

\begin{frame}{Plant structure in sun and shade}
    \centering
    \includegraphics[height=0.75\textheight]{figures/ValPea1998-3.eps}\\
    {\footnotesize Structure of sun and shade plants of \species{Heteromeles arbutifolia}
    \autocite[from][]{ValPea1998}.}
\end{frame}

\begin{frame}{Sun and shade}
    \begin{itemize}
        \item Shade plants may contain four or five times more chlorophyll per
        unit volume of chloroplast and have a higher $a/b$ ratio than
        sun plants because the light-harvesting complex increases.
        \item Both
        PS I and PS II decrease in shade, but the ratio of antenna Chl to
        reaction centre increases.
        \item Light capture capacity is increased but
        electron transfer components decrease very much in the shade.
        \item This is because under low light transport is limited by the
        low number of photons falling on the leaf, rather than by
        electron transport capacity.
    \end{itemize}
\end{frame}

\begin{frame}{Facultative and obligate shade plants}
    \begin{description}
        \item[Obligate shade plants:] Usually can grow in deep
        shade but die in full sunlight. Under high light there may
        be photodamage to the photosynthetic ``machinery''.
        \item[Facultative shade plants:] Can grow in shade (but
        usually not in very deep shade) and also in full sunlight.
        They acclimate.
        \item[Obligate sun plants:] tolerate only a little shade.
        Their capacity to acclimate is rather small.
    \end{description}
    \footnotesize This classification is rather artificial as there is a continuum
    of plant species adapted to different levels of shading, and
    also variability in the extent to which they can acclimate to
    different light levels.
\end{frame}

\begin{frame}{$A$ during a sunfleck \Discussion 5 + 2 mín}
    How does this response relate to light and carbon reactions?\\[1.5ex]
    \centering
    \includegraphics[width=0.85\textwidth]{figures/Lambers2008_22_Sunfleck_A.eps}\\
    {\footnotesize \autocite[from][]{Lambers2008}.}
\end{frame}

%%% to be moved %%%

%\begin{frame}{Soil water}
%    \vspace{-0.04\textheight}
%    \includegraphics[width=\textwidth]{figures/Crawley1.12StomataSoilWater.eps}\\
%%    \vspace{-0.02\textheight}
%    {\footnotesize Acclimation of stomata to soil water availability \autocite[from][]{Crawley1997}.}
%\end{frame}
%
%\section{Stomata}
%
%\begin{frame}{Stomata}
%    \vspace{-0.06\textheight}
%    \includegraphics[width=0.8\textwidth]{figures/IvyAbaxStomataSEM-1.eps}\\
%    \vspace{-0.02\textheight}
%    {\footnotesize Stomata on the abaxial epidermis of ivy (\species{Hedera helix}), SEM \autocite[from][]{Aphalo1991}.}
%\end{frame}
%
%\begin{frame}{Leaf cross-section}
%    \vspace{-0.06\textheight}
%    \includegraphics[width=0.8\textwidth]{figures/IvyCryofractureSEM-1.eps}\\
%    \vspace{-0.02\textheight}
%    {\footnotesize Cross-section of an ivy (muratti) leaf (\species{Hedera helix}), cryofracture SEM \autocite[from][]{Aphalo1991}.}
%\end{frame}
%
%\begin{frame}{Stoma}
%    \vspace{-0.06\textheight}
%    \includegraphics[width=0.8\textwidth]{figures/IvyStomaSEM-1.eps}\\
%    \vspace{-0.02\textheight}
%    {\footnotesize A stoma from ivy (\species{Hedera helix}), SEM \autocite[from][]{Aphalo1991}.}
%\end{frame}
%
%\begin{frame}{Stoma}
%    \vspace{-0.06\textheight}
%    \includegraphics[width=0.8\textwidth]{figures/IvyStomaFracture-1.eps}\\
%    \vspace{-0.02\textheight}
%    {\footnotesize Cross-section of a stoma from ivy  (\species{Hedera helix}), cryofracture SEM \autocite[from][]{Aphalo1991}.}
%\end{frame}
%
%\section{Canopy photosynthesis}
%
%\begin{frame}{Canopy photosynthesis}
%    \begin{itemize}
%        \item If we use the big leaf analogy to express
%        photosynthesis of a canopy we can compare the light
%        response of a leaf to that of a canopy.
%        \item In a canopy even when upper leaves are light saturated
%        (photosynthesizing at maximum rate ($A_{max}$) other lower
%        leaves can still increase $A$ if irradiance above the canopy
%        increases.
%        \item Consequently canopy photosynthesis ($\mathcal{A}$)
%        saturates at higher irradiances that $A$ of individual
%        leaves.
%    \end{itemize}
%\end{frame}
%
%\begin{frame}{Needle, shoot and canopy}
%    \vspace{-0.07\textheight}
%    \twocolumn[lcolwidth=0.57\linewidth,rcolwidth=0.37\linewidth]{
%    \includegraphics[height=0.75\textheight]{figures/Taiz9.11LightResponseNeedleShootCanopy.eps}}{%
%    \raggedright
%    \vspace{0.1\textheight}
%    {\footnotesize The shape of the light response and light saturation
%    depends on the scale at which we study the forest
%    \autocite[from][]{TaiZei2006}.}}
%\end{frame}
%
%\begin{frame}{LAI}
%    \begin{itemize}
%        \item Canopy photosynthesis not only depends on leaf
%        photosynthesis.
%        \item It depends on leaf area index (LAI), up to a limit the higher the LAI,
%        the more light is absorbed.
%        \item It depends also on how the leaves are displayed:
%        more or less horizontal or vertical.
%        \item LAI will depend on environmental conditions.
%    \end{itemize}
%\end{frame}
%
%\begin{frame}{Soil water affects LAI}
%    \vspace{-0.04\textheight}
%    \includegraphics[width=\textwidth]{figures/Crawley1.21LAISoilWater.eps}\\
%%    \vspace{-0.02\textheight}
%    {\footnotesize Leaf area index (LAI) and site water balance \autocite[from][]{Crawley1997}.}
%\end{frame}
%

%\section{Acclimation and adaptation to shade}
%
%\begin{frame}{Sun and shade}
%    \begin{itemize}
%        \item There are shade and sun adapted species.
%        \item In many species plants acclimate to the growing light
%        environment.
%        \item Some tree species produce very distinct sun and shade
%        leaves on different part of the crown.
%        \item When grown in controlled environments we call them
%        high light (HL) and low light (LL) plants. (same light quality.)
%        \item Natural shadelight differs in quality (colour) from
%        sunlight, and also in irradiance (quantity).
%    \end{itemize}
%\end{frame}
%
%\begin{frame}{Adaptation to shade}
%    \vspace{-0.07\textheight}
%    \twocolumn[lcolwidth=0.57\linewidth,rcolwidth=0.37\linewidth]{
%    \includegraphics[height=0.75\textheight]{figures/Taiz9.8LightResponseSunShadeSpecies.eps}}{%
%    \raggedright
%    \vspace{0.1\textheight}
%    {\footnotesize Adaptation to shade in different species (look also at the start of the curves
%    at low irradiances).
%    \autocite[from][]{TaiZei2006}.}}
%\end{frame}
%
%\begin{frame}{Acclimation to low light}
%    \vspace{-0.07\textheight}
%    \twocolumn[lcolwidth=0.57\linewidth,rcolwidth=0.37\linewidth]{
%    \includegraphics[height=0.75\textheight]{figures/Taiz9.10LightResponseSunShadeGrowth.eps}}{%
%    \raggedright
%    \vspace{0.1\textheight}
%    {\footnotesize Acclimation happens in response to growth conditions
%    \autocite[from][]{TaiZei2006}.}}
%\end{frame}
%
%\begin{frame}{Growth irradiance}
%    \vspace{-0.04\textheight}
%    \includegraphics[width=\textwidth]{figures/Crawley1.5HLLLAtriplexPhotosyn.eps}\\
%%    \vspace{-0.02\textheight}
%    {\footnotesize Acclimation to growth irradiance \autocite[from][]{Crawley1997}.}
%\end{frame}
%
%\begin{frame}{Acclimation to light}
%    \vspace{-0.04\textheight}
%    \includegraphics[width=\textwidth]{figures/Givnish1988-2.eps}\\
%%    \vspace{-0.02\textheight}
%    {\footnotesize Photosynthesis expressed on different bases: area, mass, or protein \autocite[from][]{Givnish1988}.}
%\end{frame}
%
%\begin{frame}{Growth irradiance}
%%    \vspace{-0.04\textheight}
%    \twocolumn[lcolwidth=0.57\linewidth,rcolwidth=0.37\linewidth]{
%    \includegraphics[height=0.6\textheight]{figures/HLLLplants-1.eps}}{%
%%    \vspace{-0.02\textheight}
%    \includegraphics[height=0.4\textheight]{figures/HLLLleaves-1.eps}}{%
%    \raggedright {\footnotesize Acclimation to growth irradiance in radish
%    (\species{Raphanus sativus})
%    \autocite[from][]{Lichtenthaler1985}.}}
%\end{frame}
%
%\begin{frame}{Growth irradiance}
%%    \vspace{-0.07\textheight}
%    \twocolumn[lcolwidth=0.70\linewidth,rcolwidth=0.24\linewidth]{
%    \includegraphics[height=\textheight]{figures/HLLLphotosynFuchsia-1.eps}}{%
%    \vspace{0.2\textheight}
%    \raggedright {\footnotesize Acclimation to growth irradiance in \species{Fuchsia
%    magellanica} (one clone)
%    \autocite[from][]{AphaloEtAl1991}.}}
%\end{frame}
%
%\begin{frame}{}
%    \includegraphics[width=\textwidth]{figures/Fuchsia_HL_photo.eps}\\
%    \includegraphics[width=\textwidth]{figures/Fuchsia_LL_photo.eps}
%\end{frame}
%
%\begin{frame}{Sun and shade leaves}
%%    \vspace{-0.04\textheight}
%    \twocolumn[lcolwidth=0.57\linewidth,rcolwidth=0.37\linewidth]{
%    \includegraphics[width=\textwidth]{figures/SunShadeLeaves-1.eps}}{%
%%    \vspace{-0.02\textheight}
%    \includegraphics[height=0.4\textheight]{figures/SunShadeLeaves-2.eps}}{%
%    \raggedright {\footnotesize Sun and shade leaves from different parts
%    of the crown of the same beech
%    (\species{Fagus sylvatica}) tree
%    \autocite[from][]{Lichtenthaler1985}.}}
%\end{frame}
%
%\begin{frame}{Sun and shade leaves}
%    \vspace{-0.04\textheight}
%    \includegraphics[width=0.9\textwidth]{figures/SunShadeLeaves-3.eps}\\
%    \vspace{-0.03\textheight}
%    {\footnotesize Shape can also be different in sun and shade leaves, e.g. in oak \autocite[from][]{}.}
%\end{frame}
%
%
%\begin{frame}{\species{Thermopsis}}
%    \vspace{-0.07\textheight}
%    \twocolumn[lcolwidth=0.47\linewidth,rcolwidth=0.47\linewidth]{
%    \includegraphics[height=0.70\textheight]{figures/ThermopsisPhoto-1.eps}}{%
%    \raggedright
%    \vspace{0.1\textheight}
%    {\footnotesize False lupin plants (\species{Thermopsis})}}
%    {\tiny (from \url{http://linnaeus.nrm.se/flora/di/faba/therm/thermon.html}).}
%\end{frame}
%
%\begin{frame}{Sun and shade plants}
%    \vspace{-0.04\textheight}
%    \includegraphics[width=\textwidth]{figures/SunShadeLeafThermopsis-1.eps}\\
%%    \vspace{-0.02\textheight}
%    {\footnotesize Acclimation of high and low light in \species{Thermopsis} \autocite[from][]{TaiZei2002}.}
%\end{frame}
%
%\begin{frame}{Plant structure in sun and shade}
%    \vspace{-0.07\textheight}
%    \twocolumn[lcolwidth=0.57\linewidth,rcolwidth=0.37\linewidth]{
%    \includegraphics[height=0.75\textheight]{figures/ValPea1998-3.eps}}{%
%    \raggedright
%    \vspace{0.1\textheight}
%    {\footnotesize Structure of sun and shade plants of \species{Heteromeles arbutifolia}
%    \autocite[from][]{ValPea1998}.}}
%\end{frame}

  \section*{References}
  \begin{frame}[t,allowframebreaks]
    \frametitle{References}
    \nocite{Taiz2014}
    \printbibliography
  \end{frame}


\end{document}
